\documentclass{amsart}
\usepackage{amssymb}
\usepackage{color}
\usepackage{hyperref}
\usepackage{tikz-cd}
\usepackage{rotating}
\newcommand*{\isoarrow}[1]{\arrow[#1,"\rotatebox{90}{\(\simeq\)}"
	]}

\newtheorem{theorem}{Theorem}[section]
\newtheorem{lemma}[theorem]{Lemma}
\newtheorem{proposition}[theorem]{Proposition}
\newtheorem{corollary}[theorem]{Corollary}
\newtheorem{conjecture}[theorem]{Conjecture}
\theoremstyle{definition}
\newtheorem{definition}[theorem]{Definition}
\newtheorem{example}[theorem]{Example}
\newtheorem{xca}[theorem]{Exercise}
\newtheorem{remark}[theorem]{Remark}
\numberwithin{equation}{section}
\usepackage[all]{xy}
\DeclareMathOperator{\Po}{Po}
\DeclareMathOperator{\Ost}{Ost}
\DeclareMathOperator{\Cl}{Cl}
\DeclareMathOperator{\disc}{disc}
\DeclareMathOperator{\Gal}{Gal}
\DeclareMathOperator{\Int}{Int}
\DeclareMathOperator{\Ker}{Ker}
\DeclareMathOperator{\Coker}{Coker}
\DeclareMathOperator{\Inf}{Inf}
\DeclareMathOperator{\Res}{Res}
\DeclareMathOperator{\res}{res}
\DeclareMathOperator{\ord}{ord}
\DeclareMathOperator{\Ver}{Ver}
\DeclareMathOperator{\Image}{Image}
\usepackage[all]{xy}

\begin{document}





\title{Solvable extensions of number fields ramified at only one prime are Ostrowski}



\author{Ehsan Shahoseini}
\address{Department of Mathematics, Tarbiat Modares University, 14115-134, Tehran, Iran}
\curraddr{} \email{ehsan\_shahoseini@modares.ac.ir}
\thanks{}


\author{Ali Rajaei$^{*}$}
\address{Department of Mathematics, Tarbiat Modares University, 14115-134, Tehran, Iran}
\curraddr{}
\email{alirajaei@modares.ac.ir}
\thanks{$^{*}$Corresponding author}





\subjclass[2020]{Primary 11R20, 11R29, 11R34}




\begin{abstract}
In this note, we show that, under a certain condition, solvable extensions of number fields ramifed at only one prime are Ostrowski. As a corollary, we deduce a generalization of \textit{Hilbert Theorem $94$} to cyclic extensions ramifed at one prime.
\end{abstract}

\maketitle

\section{Introduction}

The P\'olya group of a number field is defined in \cite{Cahen-Chabert's book} as follows:

\begin{definition} \label{definition, Polya group} 
	Let $L/\mathbb{Q}$ be a finite extension of number fields. The \textit{P\'olya group} of $L$ is the subgroup of $\Cl(L)$ generated by the classes of the \textit{Ostrowski ideals} defined as follows:  
	\begin{align} \label{equation, Ostrowski ideal}
		\Pi_{p^f}(L):=\prod_{\substack{\mathfrak{P}\in Max(\mathcal{O}_L) \\ N_{L/\mathbb{Q}}(\mathfrak{P})=p^f}} \mathfrak{P},
	\end{align}
	where $p$ is a prime number and $f$ is a positive integer. By convention, if $L$ has no ideal with norm $p^f$, then we put $\Pi_{p^f}(L)=\mathcal{O}_L$.
	We denote the P\'olya group of $K$ by $\Po(K)$. If $Po(K)=\{0\}$, we say that $K$ is a \textit{P\'olya field}.
\end{definition}

In \cite{Zantema}, Zantema proved the following important proposition:

\begin{proposition} \cite[Proposition $2.5$]{Zantema} \label{proposition, abelian extensions are Polya}
If $L/\mathbb{Q}$ is a finite abelian extension ramifying at only one prime, then $L$ is a P\'olya field.
\end{proposition}


The main goal of this paper is to generalize this result to the relative setting.
The P\'olya group has been relativized as follows \cite{ChabertI, MaarefparvarThesis}:

\begin{definition} \label{definition, relative Polya group} 
	Let $L/K$ be a finite extension of number fields. The \textit{relative P\'olya group} of $L$ over $K$ is the subgroup of $\Cl(L)$ generated by the classes of the \textit{relative Ostrowski ideals} defined as follows: 
	\begin{align} \label{equation, relative Ostrowski ideal}
		\Pi_{\mathfrak{p}^f}(L/K):=\prod_{\substack{\mathfrak{P}\in Max(\mathcal{O}_L) \\ N_{L/K}(\mathfrak{P})=\mathfrak{p}^f}} \mathfrak{P},
	\end{align}
	where $\mathfrak{p} \in Max(\mathcal{O}_K)$ and $f$ is a positive integer. By convention, if $L$ has no ideal with norm $\mathfrak{p}^f$ (over $K$), then we put $\Pi_{\mathfrak{P}^f}(L/K)=\mathcal{O}_L$.
	We denote the relative P\'olya group of $L$ over $K$ by $\Po(L/K)$. In particular, $\Po(L/\mathbb{Q})=\Po(L)$ and $\Po(L/L)=\Cl(L)$.
\end{definition}

\begin{remark} \label{remark, elements of relative Polya group}
Let $L/K$ be a finite Galois extension of number fields with Galois group $G$, then $Po(L/K)=\frac{I(L)^G}{P(L)^G}$. Hence, an ideal calss $[\mathfrak{a}] \in Cl(L)$ is contained in $Po(L/K)$ if and only if there exists an ideal $\mathfrak{a}_1$ such that $\mathfrak{a}_1 \in [\mathfrak{a}]$ and $\mathfrak{a}_1^{\sigma} = \mathfrak{a}_1$ for all $\sigma \in G$.
\end{remark}

Then, in \cite{Ehsan Thesis} (and \cite{SRM}) a modification of the relative Polya group has been introduced as follows:

\begin{definition} \label{definition, Ostrowski quotient}
Let $L/K$ be a finite extension of number fields. Then, the \textit{Ostrowski quotient} of $L$ over $K$, $\Ost(L/K)$, is defined as follows:
	\begin{equation*} \label{equation, Ostrowski quotient}
	\Ost(L/K) := \frac{\Po(L/K)}{\Po(L/K) \cap \epsilon_{L/K}(\Cl(K))},
	\end{equation*} 
	where $\epsilon_{L/K}$ is the \textit{Capitulation map} from $\Cl(K)$ to $\Cl(L)$.
	In particular, $\Ost(L/\mathbb{Q})=\Po(L/\mathbb{Q})=\Po(L)$ and $\Ost(L/L)=\{0\}$. Note that if $L/K$ is Galois, then we have $\epsilon_{L/K}(\Cl(K)) \subseteq \Po(L/K)$ and thus:
	\begin{equation*} \label{equation, Ostrowski quotient Galois case}
	\Ost(L/K) := \frac{\Po(L/K)}{\epsilon_{L/K}(\Cl(K))}.
	\end{equation*} 
If $Ost(L/K)=\{0\}$, we say that $L/K$ is an \textit{Ostrowski extension}.
\end{definition}

In \cite{SRM} the following conditional generalization of Proposition \ref{proposition, abelian extensions are Polya} has been proved:

\begin{theorem} \cite[Theorem $3.6$]{SRM} \label{theorem, abelian extensions are Ostrowski, conditional}
Let $K/F$ be a finite abelian extension of number fields such that only one prime of $F$ is ramified in $K$. Let $L$ be the ray class field of $F$ for the modulus $\mathfrak{c}(K/F)$, where $\mathfrak{c}(K/F)$ denotes the conductor of $K$ over $F$. If $L/F$ is Ostrowski, then so is $K/F$.
\end{theorem} 

In the next section we will drop the condition of Ostrowskiness of ray class field $L$ of $F$ in Theorem \ref{theorem, abelian extensions are Ostrowski, conditional} and we will go from abelian extensions to solvable extensions, but we will need to impose a certain condintion on our extension.

In \cite[p.$163$]{Zantema}, Zantema proved the following theorem:

\begin{theorem} \label{theorem, Zantema's exact sequence}
Let $L$ be a number field such that $L/\mathbb{Q}$ is a finite Galois extension with Galois group $G$. Then, we have the following exact sequence:
\begin{equation*} \label{equation, Zantema's exact sequence}
0 \rightarrow H^1(G,U_L) \rightarrow \bigoplus_{p \: prime} \frac{\mathbb{Z}}{e_{p(L/\mathbb{Q})}\mathbb{Z}} \rightarrow \Po(L) \rightarrow 0,
\end{equation*}
where $U_L$ is the unit group of $L$ and $e_p(L/\mathbb{Q})$ is the ramification index of $p$ in $L$.
\end{theorem}

By use of some Galois cohomology results of Brumer-Rosen \cite{Brumer-Rosen}, the above exact sequence of Zantema generalizes as follows:

\begin{theorem} \cite[Theorem 2.2]{MR2} \label{theorem, generalization of the Zantema's exact sequence}
	Let $L/K$ be a finite Galois extension of number fields with Galois group $G$. Then the following sequence is exact:
	{\footnotesize
		\begin{equation*} \label{equation, BRZ exact sequence}
				\tag{BRZ} \qquad 
			0 \rightarrow \Ker({\epsilon}_{L/K}) \rightarrow H^1(G,U_L) \rightarrow \bigoplus_{\mathfrak{p} \in Max(\mathcal{O}_K)} \frac{\mathbb{Z}}{e_{\mathfrak{p}(L/K)}\mathbb{Z}} \rightarrow \frac{\Po(L/K)}{{\epsilon}_{L/K}(\Cl(K))}=\Ost(L/K) \rightarrow 0.
	\end{equation*}}
\end{theorem}





\section{Main Results}

In this section, we state and prove our main result and a corollary of it which is a generalization of \textit{Hilbert Theorem $94$} to cyclic extensions ramified at one prime. First, we prove a lemma which is essential for the main theorem.

\begin{lemma} \label{lemma, short exact sequence of Ostrowski qoutients}
Let $F \subseteq K \subseteq L$ be a tower of finite extensions of number fields such that $L/F$ and $K/F$ are Galois. 
Then the following sequence is exact:
\begin{equation} \label{equation, short exact sequences for Ostrowski qoutients}
\Ost(K/F) \xrightarrow{\psi} \Ost(L/F) \xrightarrow{\varphi} \Ost(L/K),
\end{equation}
where 
\begin{align*}
\psi: \Ost(K/F) &\rightarrow \Ost(L/F) \\
[\mathfrak{a}]  \left( \mathrm{mod}\, \, \epsilon_{K/F}(\Cl(F)) \right) & \mapsto \epsilon_{L/K}([\mathfrak{a}])  \left( \mathrm{mod}\, \, \epsilon_{L/F}(\Cl(F)) \right),
\end{align*} 
and
\begin{align*}
\varphi: \Ost(L/F) &\rightarrow \Ost(L/K)  \\
[\mathfrak{b}]  \left( \mathrm{mod}\, \, \epsilon_{L/F}(\Cl(F)) \right) & \mapsto [\mathfrak{b}]  \left( \mathrm{mod}\, \, \epsilon_{L/K}(\Cl(K)) \right).\nonumber
\end{align*}
\end{lemma}

\begin{proof}
Since $K/F$ and $L/F$ are Galois extensions, one has
\begin{equation*}
\epsilon_{L/K}(\Po(K/F)) \subseteq \Po(L/F) \subseteq \Po(L/K),
\end{equation*}
see \cite[Lemma 2.10]{MR2}.  In \cite[proof of Theorem 3.14]{SRM} it is shown that $\psi$ is well-defined. Likewise, one can  show that $\varphi$ is also well-defined:

For $[\mathfrak{b}_1], [\mathfrak{b}_2] \in \Po(L/F)$, if $[\mathfrak{b}_1]=[\mathfrak{b}_2] \left( \mathrm{mod}\, \, \epsilon_{L/F}(\Cl(F)) \right) $, then $[\mathfrak{b}_1].[\mathfrak{b}_2]^{-1}=\epsilon_{L/F}([\mathfrak{c}])$ for some $[\mathfrak{c}] \in \Cl(F)$. Hence:
\begin{equation*}
[\mathfrak{b}_1].[\mathfrak{b}_2]^{-1}=\epsilon_{L/K}\left(\epsilon_{K/F}([\mathfrak{c}])\right)  \in \epsilon_{L/K}(\Cl(K)).
\end{equation*}


Now, by definition of $\varphi$ and $\psi$, we have:
\begin{equation*}
\varphi \circ \psi  \left([\mathfrak{a}]\right) \in \epsilon_{L/K}(\Cl(K)) \quad \forall \, [\mathfrak{a}] \in \Ost(K/F),
\end{equation*}
i.e., $\Image(\psi) \subseteq \Ker(\varphi)$. To check the reverse inclusion, let $[\mathfrak{b}] \in \Po(L/F)$ be such that its class in $\Ost(L/F)$ is contained in $\Ker(\varphi)$. Then, there exists $[\mathfrak{c}] \in \Cl(K)$ for which $\epsilon_{L/K}([\mathfrak{c}])=[\mathfrak{b}]$. We want to show that $[\mathfrak{c}] \in \Po(K/F)$. Since $[\mathfrak{b}] \in \Po(L/F)$, by Remark \ref{remark, elements of relative Polya group} there exists $\mathfrak{b}_1 \in [\mathfrak{b}]$ such that $\mathfrak{b}_1^{\sigma} = \mathfrak{b}_1$ for all $\sigma \in \Gal(L/F)$. Put $\mathfrak{c}_1:=\mathfrak{b}_1 \cap \mathcal{O}_K$. Then, we have $\mathfrak{c}_1 \in [\mathfrak{c}]$ and $\mathfrak{c}_1^{\delta} = \mathfrak{c}_1$ for all $\delta \in \Gal(K/F)$. Thus, again by Remark \ref{remark, elements of relative Polya group}, $\mathfrak{c} \in \Po(K/F)$ and we have:
\begin{equation*}
\psi\left([\mathfrak{c}] \, \, \mathrm{mod}\, \epsilon_{K/F}(\Cl(F)) \right)=[\mathfrak{b}] \, \, \mathrm{mod}\, \epsilon_{L/F}(\Cl(F)),
\end{equation*}
i.e., $[\mathfrak{b}]  \left(\mathrm{mod}\, \epsilon_{L/F}(\Cl(F))\right) \in \Image(\psi)$ and the proof is complete.
\end{proof}


Let $L/K$ be a finite Galois extension with Galois group $G$. We want to describe the map $\lambda: H^1(G,U_L) \rightarrow \bigoplus_{\mathfrak{p} \in Max(\mathcal{O}_K)} \frac{\mathbb{Z}}{e_{\mathfrak{p}(L/K)}\mathbb{Z}}$ in Theorem \ref{theorem, generalization of the Zantema's exact sequence} (as in \cite[p.80]{Aviles}). For 
its description, we need to review some preliminary facts. Firstly, we have $H^1(G,U_L) \simeq \frac{P(L)^G}{P(K)}$; see \cite[p.370]{MR2}. Note that the projection $P(L)^G=(L^\times/U_L)^G \twoheadrightarrow H^1(G,U_L)$ 
maps a class $[\beta U_L] \in (L^\times/U_L)^G$ to the cohomology class $[\xi]$ which is represented by the $1$-cocycle $\xi_{\beta}:G \rightarrow U_L$ defined by $\xi_{\beta}(\sigma)=\frac{\beta^{\sigma}}{\beta}$ for all $\sigma \in G$; see \cite[p.78]{Aviles}.
 Also, for any ramified prime $\mathfrak{p}$ of $K$ in this extension, we have $\frac{\mathbb{Z}}{e_{\mathfrak{p}(L/K)}\mathbb{Z}} \simeq H^1(G_{\mathfrak{P}}, U_{L_{\mathfrak{P}}})$, for $\mathfrak{P}$ a fixed prime of $L$ dividing $\mathfrak{p}$ and $G_{\mathfrak{P}}$ the decomposition group at $\mathfrak{P}$ ($L_{\mathfrak{P}}$ is 
 the completion of $L$ at $\mathfrak{P}$); see \cite[Proof of Lemma 2.3]{Aviles}. Now, Let $R$ be the set of ramified primes of the extension $L/K$ and consider the map $\lambda: H^1(G,U_L) \rightarrow \bigoplus_{\mathfrak{p} \in Max(\mathcal{O}_K)} \frac{\mathbb{Z}}{e_{\mathfrak{p}(L/K)}\mathbb{Z}} \simeq \bigoplus_{\mathfrak{p} \in R} H^1(G_{\mathfrak{P}}, U_{L_{\mathfrak{P}}})$. Then, 
 the map $\lambda$ can be described as follows: let $[c] \in H^1(G,U_L)$ be represented by the $1$-cocycle $\xi:G \rightarrow U_L$ with $\xi(\sigma)=\frac{\beta^{\sigma}}{\beta}$, where $[\beta U_L] \in (L^\times/U_L)^G$. Then, for $\mathfrak{p} \in R$, the $\mathfrak{p}$-component of $\lambda(c)$ is the cohomology class in $H^1(G_{\mathfrak{P}}, U_{L_{\mathfrak{P}}})$ represented by 
 the $1$-cocycle $\xi_{\mathfrak{p}}:G_{\mathfrak{P}} \rightarrow U_{L_{\mathfrak{P}}}$ given by $\xi_{\mathfrak{P}}(\delta)=\frac{\beta^{\delta}}{\beta}$ for $\delta \in G_{\mathfrak{P}}$; see \cite[p.80]{Aviles}.

Now, we are ready to prove the main theorem in three steps.

\begin{proposition} \label{proposition, cyclic of prime degree case}
Let $L/K$ be a finite cyclic extension of prime degree $p$ with Galois group $G$ ramified at only one prime. Then, $L/K$ is Ostrowski.
\end{proposition}

\begin{proof}
Let $L/K$ be ramified at the prime $\mathfrak{p}$ with ramification index $p$.The third term of the BRZ exact sequence \eqref{equation, BRZ exact sequence} is 
equal to $\frac{\mathbb{Z}}{p\mathbb{Z}}$, which by remarks after Lemma \ref{lemma, short exact sequence of Ostrowski qoutients} is equal to $ H^1(G_{\mathfrak{P}},U_{L_{\mathfrak{P}}})$, where $\mathfrak{P}$ is a fixed prime of $L$ dividing $\mathfrak{p}$. By the explicit description of the map $\lambda: H^1(G,U_L) \rightarrow \frac{\mathbb{Z}}{p\mathbb{Z}} \simeq  H^1(G_{\mathfrak{P}},U_{L_{\mathfrak{P}}})$, we see that $\lambda$ is not the zero map. Now, since the order of the target of the map $\lambda$ is a prime number, $\lambda$ is surjective. Hence, by the BRZ exact sequence \eqref{equation, BRZ exact sequence} we get that $\Ost(L/K)=\{0\}$, i.e. $L/K$ is Ostrowski.
\end{proof}

\begin{proposition} \label{proposition, totally ramified case}
Let $L/K$ be a finite solvable extension of number fields totally ramified at its unique ramifed prime $\mathfrak{p}$. Then, $L/K$ is Ostrowski.
\end{proposition}

\begin{proof}
Since any solvable extension can be written as a tower of abelian extensions and any abelian extension can be written as a tower of cyclic extensions of prime degree, $L/K$ can be written as a tower of cyclic extensions $K=F_0 \subseteq F_1 \subseteq \dots F_{n-1} \subseteq F_n=L$, where each $F_i/F_{i-1}$ is cyclic of prime degree $p_i$. As $L/K$ is totally ramified at $\mathfrak{p}$, we get that all $F_i/F_{i-1}$ are totally ramified at the unique prime of $F_{i-1}$ above $\mathfrak{p}$, so by Proposition \ref{proposition, cyclic of prime degree case} all extensions $F_i/F_{i-1}$ are Ostrowski. Now, consider the tower $F_0 \subseteq F_1 \subseteq F_2$. By Lemma \ref{lemma, short exact sequence of Ostrowski qoutients}, we have that the sequence:
\begin{equation*}
\Ost(F_1/F_0) \rightarrow \Ost(F_2/F_0) \rightarrow \Ost(F_2/F_1)
\end{equation*}
is exact. But, the left term and the right term are trivial, hence the middle one is trivial, too. Then, consider the tower $F_0 \subseteq F_2 \subseteq F_3$. Again by Lemma \ref{lemma, short exact sequence of Ostrowski qoutients}, we get the exact sequence:
\begin{equation*}
\Ost(F_2/F_0) \rightarrow \Ost(F_3/F_0) \rightarrow \Ost(F_3/F_2).
\end{equation*}
Both left and right terms are trivial agian, hence the middle one is, too. Using induction and considering the tower $F_0 \subseteq F_{n-1} \subseteq F_n$, we get that $\Ost(F_n/F_0)=\Ost(L/K)$ is trivial.
\end{proof}

\begin{remark}
In Proposition \ref{proposition, totally ramified case}, since the extension $L/K$ (with Galois group $G$) has a totally ramified prime, the map $\epsilon_{L/K}$ is injective (\cite[Theorem 2.1]{Masley}) and thus by the BRZ exact sequence \eqref{equation, BRZ exact sequence} we find that $H^1(G,U_L) \simeq \frac{\mathbb{Z}}{[L:K]\mathbb{Z}}$.
\end{remark}

\begin{theorem} \label{theorem, solvable extensions are Ostrowski, conditional}
Let $L/K$ be a solvable extension of number fields with Galois group $G$ ramified at only one prime $\mathfrak{p}$. Also, let $I_{\mathfrak{P}}$ be the inertia at $\mathfrak{P}$, where $\mathfrak{P}$ is a fixed prime of $L$ dividing $\mathfrak{p}$, and let $\mathfrak{p}$ remain prime in $L^{I_{\mathfrak{P}}}/K$. Then, $L/K$ is Ostrowski.
\end{theorem}

\begin{proof}
Since $\mathfrak{p}$ remains prime in $L^{I_{\mathfrak{P}}}/K$ and is totally ramified in $L/L^{I_{\mathfrak{P}}}$, the prime $\mathfrak{P}$ in $L$ that divides $\mathfrak{p}$ is unique and also $G \simeq D_{\mathfrak{P}}$, where $D_{\mathfrak{P}}$ is the decomposition group at $\mathfrak{P}$. Since $I_{\mathfrak{P}}$ is a normal subgroup of $D_{\mathfrak{P}}\simeq G$, by Lemma \ref{lemma, short exact sequence of Ostrowski qoutients} we get the following exact sequecne for the tower $K \subseteq L^{I_{\mathfrak{P}}} \subseteq L$:
\begin{align*}
\Ost(L^{I_{\mathfrak{P}}}/K) \rightarrow \Ost(L/K) \rightarrow \Ost(L/L^{I_{\mathfrak{P}}}).
\end{align*}
Since $L^{I_{\mathfrak{P}}}/K$ is unramified, the BRZ exact sequence \eqref{equation, BRZ exact sequence} shows that $\Ost(L^{I_{\mathfrak{P}}}/K)$ is trivial. Also, since $L/L^{I_{\mathfrak{P}}}$ is solvable and is totally ramified at only one prime $\mathfrak{p}$, by Proposition \ref{proposition, totally ramified case} it is Ostrowski. Thus $L/K$ is Ostrowski.
\end{proof}

\begin{remark}
Note that in the Theorem \ref{theorem, solvable extensions are Ostrowski, conditional}, we can not drop the condition ''$\mathfrak{p}$ remains prime in $L^{I_{\mathfrak{P}}}/K$``. Because in non-abelian extensions the inertia group at a prime can be a non-normal subgroup of the Galois group; it is just a normal subgroup of the decomposition group at that prime.
\end{remark}

Hilbert Theorem $94$ says that for an unramified cyclic extension $L/K$, $[L:K] \mid \#\ker(\epsilon_{L/K}(Cl(K)))$. In fact, in \cite{SRM}, using BRZ exact sequence \eqref{equation, BRZ exact sequence} we showed that $\#\ker(\epsilon_{L/K}(\Cl(K)))=[L:K] \# \hat{H}^0(G,U_L)$, where $G$ is the Galois group of $L/K$. Note that in this case, for any prime $\mathfrak{p}$ of $K$ we have $[L:K]=f_{\mathfrak{p}}g_{\mathfrak{p}}$, where $f_{\mathfrak{p}}$ is the residue degree at $\mathfrak{p}$ and  $g_{\mathfrak{p}}$ is the number of primes of $L$ dividing $\mathfrak{p}$. As a corollary of Theorem \ref{theorem, solvable extensions are Ostrowski, conditional}, we can prove the following generalization of Hilbert Theorem $94$ for a ramified case:

\begin{corollary}
Let $L/K$ be a cyclic extension of number fields with Galois group $G$ ramified at only one finite prime $\mathfrak{p}$ and unramified at infinite primes. Assume there exists a unique prime $\mathfrak{P}$ of $L$ which divides $\mathfrak{p}$. Also, let $e_{\mathfrak{p}}$ and $f_{\mathfrak{p}}$ denote the ramification index and the residue degree of $\mathfrak{p}$, respectively. Then, $f_{\mathfrak{p}}  \#\hat{H}^0(G,U_L) \mid \#\ker(\epsilon_{L/K}(\Cl(K)))$ (note that $f_{\mathfrak{p}} =\frac{[L:K]}{e_{\mathfrak{p}}}$).
\end{corollary}
 
\begin{proof}
Since the extension $L/K$ satisfies the conditions of Theorem \ref{theorem, solvable extensions are Ostrowski, conditional}, $\Ost(L/K)=\{0\}$. Thus, by the BRZ exact sequence, we get $\#\ker(\epsilon_{L/K}(\Cl(K)))=\frac{\#H^1(G,U_L)}{e_{\mathfrak{p}}}$. As $L/K$ is cyclic, we can use the \textit{Herbrand quotient}:
\begin{align*} \label{equation, Herbrand quotients of L over K}
Q(G,U_L)=\frac{\# \hat{H}^0(G, U_L)}{\#{H}^1(G,U_L)}.
\end{align*}
Also, we have:
\begin{align*}
Q(G,U_L)= \frac{2^s}{[L:K]},
\end{align*}
where $s$ is the number of infinite primes of $K$ ramified in $L$, \cite[IX,\S4,Corollary 2]{Lang}. By assumption, $L/K$ is unramified at infinite primes which implies that $s=0$ i.e. $Q(G,U_L)=\frac{1}{[L:K]}$. So, $[L:K] \#\hat{H}^0(G,U_L) \mid \#H^1(G,U_L)$. Hence, we get that $\frac{[L:K] \#\hat{H}^0(G,U_L)}{e_{\mathfrak{p}}} = f_{\mathfrak{p}} \#\hat{H}^0(G,U_L) \mid \#\ker(\epsilon_{L/K}(\Cl(K)))$.
\end{proof}



\bibliographystyle{amsplain}

\begin{thebibliography}{20}
	
	\baselineskip=.5cm
	
	
	\bibitem{Brumer-Rosen} A. Brumer, M. Rosen, \textit{Class number and ramification in number fields.} Nagoya Math. J. \textbf{23} (1963), 97--101.
	
	
	\bibitem{Cahen-Chabert's book} P. J. Cahen, J. L. Chabert \textit{Integer-valued polynomials.} Mathematical Surveys and Monographs \textbf{48}, Amer. Math. Soc., Providence, 1997.
	
	
	\bibitem{ChabertI} J. L. Chabert, \textit{From P\'olya fields to P\'olya groups (I) Galois extensions.} J. Number Theory \textbf{203} (2019), 360--375.
	
	
	
\bibitem{Aviles} C. D. Gonz\'alez-Avil\'es,  \textit{Capitulation, ambiguous classes and the cohomology of the units}. J. Reine Angew. Math. \textbf{613} (2007), 75--97.
	

\bibitem{Lang} S. Lang, \textit{Algebraic Number Theory}. Springer, New York, 1986.
	
	
	\bibitem{MaarefparvarThesis} A. Maarefparvar, \textit{P\'olya fields of small degree and low ramification.} PhD thesis (in Persian), Tarbiat Modares University, December 2017.

	
	
	\bibitem{MR2} A. Maarefparvar, A. Rajaei, \textit{Relative P\'olya group and P\'olya dihedral extensions of $\mathbb{Q}$}. J. Number Theory \textbf{207} (2020), 367--384.
	
	
	\bibitem{Masley} J. M. Masley, \textit{Class numbers of real cyclic number fields with small conductor.} Compositio Math. \textbf{37} (1978), no. 3, 297--319.
	
	
	
	\bibitem{Ehsan Thesis} E. Shahoseini, \textit{On Class groups and Brumer-Rosen's exact sequence.} PhD thesis (in Persian), Tarbiat Modares University, June 2021.
	
\bibitem{SRM} E. Shahoseini, A. Rajaei, A. Maarefparvar, \textit{Ostrowski quotients for finite extensions of number fields.} Pacific Journal of Math. (2022), \textbf{321} (2022), No. 2, 415--429.
	
	

	
	
	\bibitem{Zantema} H. Zantema, \textit{Integer valued polynomials over a number field.} Manuscripta Math. \textbf{40} (1982), 155--203. 
	
\end{thebibliography}



\end{document}
